This section defines the principal system-level requirements for the \abbr{EWAS}. The requirements are grouped by functional area and consolidated in Table~\ref{tab:system-requirements} to support review, traceability, and verification planning.

\begingroup
\footnotesize
\setlength{\parskip}{0pt}
\setlength{\LTpre}{0.5\baselineskip}
\setlength{\LTpost}{0pt}
\renewcommand{\arraystretch}{1.18}
\setlength{\tabcolsep}{5pt}
\rowcolors{2}{gray!8}{white}

\begin{longtable}{@{}>{\raggedright\arraybackslash\bfseries}p{0.11\textwidth}>{\raggedright\arraybackslash}p{0.23\textwidth}>{\raggedright\arraybackslash}p{0.58\textwidth}@{}}
\caption{EWAS system requirements}\label{tab:system-requirements}\\
\toprule
\rowcolor{gray!25}
\textbf{ID} & \textbf{Requirement} & \textbf{Statement} \\
\midrule
\endfirsthead

\multicolumn{3}{c}{\tablename~\thetable\ continued from the previous page}\\
\toprule
\rowcolor{gray!25}
\textbf{ID} & \textbf{Requirement} & \textbf{Statement} \\
\midrule
\endhead

\midrule
\multicolumn{3}{r}{Continued on the next page}\\
\endfoot

\bottomrule
\endlastfoot

\rowcolor{blue!12}
\multicolumn{3}{@{}l}{\textbf{Top-Level System Requirements}}\\
R-001 & Non-Offensive Operation & The system shall not provide targeting information for weapons systems, weapon guidance or fire-control data, interceptor command-and-control functions, or direct control of military response assets. \\
R-002 & Scope Limitation & The system shall not be required to provide definitive missile-type identification, final warhead characterization, exact explosive-yield estimation, or comprehensive post-impact damage assessment. \\
R-003 & Authorized User Access & System outputs shall be accessible only to authorized governmental and institutional users, including crisis-management and civil-protection authorities. \\
R-004 & Data Usage & Historical mission data shall be used solely for algorithm improvement, system calibration, performance verification, operator training, and post-event technical analysis. \\

\rowcolor{blue!12}
\multicolumn{3}{@{}l}{\textbf{Constellation and Coverage Requirements}}\\
R-005 & Space Segment Architecture & The \abbr{EWAS} space segment shall consist of six \abbr{GEO} satellites and 72 \abbr{LEO} satellites operating at an altitude of 1,200~km. \\
R-006 & Continuous Coverage & The constellation shall provide continuous surveillance of the designated monitoring region on a 24-hour basis. \\
R-007 & Boost-Phase Detection Coverage & The orbital architecture shall provide observation capability during the missile boost phase to enable the earliest feasible launch detection. \\
R-008 & Coverage Availability & The constellation shall maintain continuous temporal coverage of the designated surveillance region without operational gaps that would prevent launch detection. \\

\rowcolor{blue!12}
\multicolumn{3}{@{}l}{\textbf{Performance Requirements}}\\
R-009 & Launch Detection & The system shall detect ballistic missile launch events by observing characteristic infrared signatures generated by missile propulsion systems. \\
R-010 & Probability of Detection & The probability of detection shall exceed 95\% for missile threats defined within the mission scope. \\
R-011 & End-to-End Latency & The elapsed time from the first \abbr{GEO} detection of a candidate launch event to delivery of the preliminary warning product at an authorized national authority interface shall not exceed 60 seconds. The warning shall follow processing of observations from the two selected \abbr{LEO} satellites. \\
R-012 & Preliminary Impact Estimation & Following launch detection, the ground segment shall generate an initial impact-area estimate using the preliminary trajectory and telemetry downlinked by the two tracking \abbr{LEO} satellites. \\
R-013 & Uncertainty Quantification & Impact-area estimates shall be presented as probabilistic confidence regions with clearly defined confidence levels and uncertainty boundaries. \\
R-014 & Dynamic Update Capability & The ground segment shall update the refined trajectory and impact-area estimates as additional tracking data become available. \\
R-015 & Operational Availability & The overall system availability shall exceed 99.9\%. \\
R-017 & System Redundancy & The system architecture shall incorporate sufficient redundancy to maintain operational performance during component failures. \\
R-019 & Event Classification & The system shall automatically distinguish missile launch events from non-threat infrared sources through onboard and ground-based processing. \\

\rowcolor{blue!12}
\multicolumn{3}{@{}l}{\textbf{Space Segment Requirements}}\\
R-020 & Infrared Sensor Payload & Each satellite shall carry infrared sensing equipment capable of detecting missile-plume signatures against background environmental conditions. \\
R-021 & Sensor Coverage & The payload shall support wide-area surveillance while maintaining sufficient sensitivity to satisfy launch-detection requirements. \\
R-022 & Onboard Event Processing & Satellites shall perform onboard event filtering and preliminary event characterization to reduce communication latency and bandwidth requirements. \\
R-023 & Preliminary Trajectory Estimation & The onboard processing subsystems of the two cued \abbr{LEO} satellites shall use stereoscopic observations to calculate a preliminary trajectory estimate before transmitting it to the ground segment. \\
R-024 & Pointing Stability & The spacecraft shall maintain pointing accuracy and stability sufficient to achieve mission detection-performance requirements. \\
R-025 & Stray-Light Mitigation & The spacecraft design shall minimize stray-light contamination that could degrade sensor performance. \\

\rowcolor{blue!12}
\multicolumn{3}{@{}l}{\textbf{Ground Segment Requirements}}\\
R-026 & Data Fusion & The ground segment shall support real-time fusion of all relevant available \abbr{GEO} and \abbr{LEO} observations associated with a detected event. \\
R-027 & Threat Assessment Processing & The ground segment shall fuse the available \abbr{LEO} observations and preliminary trajectory estimates to calculate a refined trajectory estimate and probabilistic impact-area prediction in near-real time. \\
R-028 & Decision Support Interface & The user interface shall visualize the launch location, estimated trajectory, predicted impact area, confidence level, and uncertainty region. \\
R-029 & Human Factors Design & The interface shall be designed to minimize operator misinterpretation during high-stress situations. \\
R-030 & Command and Control & The ground segment shall provide redundant command-and-control capabilities for constellation management. \\

\rowcolor{blue!12}
\multicolumn{3}{@{}l}{\textbf{Communications Requirements}}\\
R-031 & Secure Communications & All communications shall employ encryption mechanisms compliant with applicable national security requirements. \\
R-032 & Anti-Jamming Protection & Communication links shall incorporate protection against jamming, spoofing, and unauthorized access. \\
R-033 & Near-Real-Time Data Transmission & The communication subsystem shall support near-real-time transmission of sensor and telemetry data between the space and ground segments. \\

\rowcolor{blue!12}
\multicolumn{3}{@{}l}{\textbf{Data Management and Security Requirements}}\\
R-034 & Role-Based Access Control & The system shall implement Role-Based Access Control (\abbr{RBAC}) to restrict access to authorized users. \\
R-035 & Secure Data Storage & All archived mission data shall be stored in encrypted repositories. \\
R-036 & Data Retention & Historical mission data shall be retained for a minimum period of five years. \\
R-037 & Audit Logging & The system shall maintain audit records of user actions, system events, and warning-dissemination activities. \\

\rowcolor{blue!12}
\multicolumn{3}{@{}l}{\textbf{Verification Requirements}}\\
R-038 & Simulation-Based Verification & System performance shall be verified through representative simulation campaigns covering a range of missile-launch scenarios and environmental conditions. \\
R-039 & Hardware-in-the-Loop Testing & Critical sensing and processing functions shall undergo Hardware-in-the-Loop (\abbr{HIL}) testing prior to operational deployment. \\
R-040 & Operational Validation & End-to-end operational exercises shall demonstrate compliance with detection, latency, availability, and warning-dissemination requirements. \\
R-041 & Performance Monitoring & Archived operational data shall be periodically analyzed to support algorithm refinement and continuous system improvement. \\

\rowcolor{blue!12}
\multicolumn{3}{@{}l}{\textbf{Additional Operational Requirements (proposed change)}}\\
R-042 & GEO Event Reporting & On detecting a candidate thermal event, the detecting \abbr{GEO} satellite shall transmit event time, estimated location and available characterization data to the ground segment. \\
R-043 & Ground-Directed LEO Selection & Following receipt of a \abbr{GEO} event report, the ground segment shall select two suitable \abbr{LEO} tracking satellites taking account of event location, observation geometry and link availability. \\
R-044 & Tracking Command Routing & The ground segment shall transmit tracking requests to the selected \abbr{LEO} satellites through available supported ground-station links or, where available, \abbr{ISL} relay paths. \\
R-045 & Post-LEO Preliminary Warning & The Early Warning System shall automatically issue the first preliminary warning only after the ground segment has received and assessed observations from the two selected \abbr{LEO} satellites and produced a refined trajectory and impact-area estimate. \\
R-046 & Operator Review and Revision & The Mission Control Center shall present each issued preliminary warning to an authorized operator and support subsequent confirmation, revision or cancellation without requiring approval before the first dispatch. \\
R-047 & Record Retention of Events & The Data Archiving System shall retain the applicable event data, processed results, warning versions, delivery acknowledgements and audit logs irrespective of whether the operator later confirms or cancels an event, subject to R-004 and R-036. \\
R-048 & Boundary of Warning Delivery & The system shall deliver warning products to authorized national authority interfaces; subsequent public-warning activation shall remain outside the \abbr{EWAS} system boundary. \\
R-049 & Protected Region Mission Focus & The system shall assess whether launch events, including those originating outside the Baltic region, may threaten the Baltic region, and shall prioritize warning products accordingly. \\

\end{longtable}
\endgroup

\paragraph{Requirement configuration status.}
R-016 (single-failure fault-tolerance guarantee) and R-018 (numerical false-alarm-rate limit) were withdrawn following the project team's agreed baseline decisions and are intentionally not renumbered or reproduced in the active requirement table. R-015, R-017 and R-030 remain applicable; removing R-016 does not remove all system redundancy obligations. Requirements R-042--R-049 are newly proposed derived operational requirements included to capture the revised two-segment, ground-directed tracking and post-LEO warning concept; they require team review and formal baseline acceptance. The 60-second requirement in R-011 is a design target, not a verified achievement, and remains a key open issue pending communication and end-to-end latency analysis.
